Open-set recognition (OSR) addresses classification settings in which a model must preserve recognition of known classes while identifying samples that do not belong to any known class. This problem is especially important in cyber and electromagnetic sensing, where the operational value of a detector depends not only on assigning labels to familiar behaviors but also on refusing overconfident assignments for novel or adversarial behaviors. Confidence-based OSR methods provide a natural starting point because they operate on the decision structure of a closed-set classifier rather than requiring a separate generative model for every possible unknown.

The closest methodological predecessor to this work is varMax. Baye et al. introduced varMax as a confidence-based zero-day attack recognition method that uses maximum-score and variance-style evidence from classifier outputs to reject unknown samples~\cite{baye2024varmax}. Broggi et al. further develop the uncertainty and novelty-management perspective behind varMax, emphasizing the use of classifier confidence structure for open-set decisions~\cite{broggi2025varmax}. Our work inherits this view that novelty can be detected from the geometry of closed-set confidence outputs, but differs in two ways. First, DQNGuard conditions calibration on the predicted class rather than applying only a global rejection rule. Second, it uses a known-budget threshold so that unknown detection is constrained by an explicit limit on known-sample rejection.

A second line of related work uses reinforcement-learning-inspired decision layers for open-set intrusion detection. Tiwari et al. propose DQN-IDS, a deep reinforcement learning approach for open-set-enabled intrusion detection~\cite{tiwaridqn}. That work motivates treating OSR as a decision process over confidence states rather than as a pure feature-extraction problem. In this paper, we do not claim deployment-time continual reinforcement learning. Instead, we evaluate a DQN-IDS-style confidence head as a comparison point and use DQNGuard as a calibrated decision layer over fixed backbone outputs. This separation is useful experimentally because performance differences can be attributed to the OSR decision rule rather than to a different feature extractor or online learning policy.

Open-set RF and electromagnetic signal analysis has also explored multi-domain signal representations. Wei et al. exploit raw IQ, FFT, and DCT domains for detecting unclassified electromagnetic signals, using domain-specific autoencoders and feature fusion before open-set discrimination~\cite{wei2024multidomain}. Trott et al. similarly study RF modulation classifiers in the presence of novel samples and motivate multi-domain representations for RF OSR~\cite{trott2026model}. The present work adopts the same broad principle that RF novelty may appear differently across time-domain, frequency-domain, and transformed representations. However, rather than fusing autoencoder latents, our backbone processes IQ, FFT, DCT, and polar components through convolutional feature paths and then exposes logits, softmax probabilities, predicted labels, and intermediate features to the OSR decision layer.

This paper is therefore positioned at the intersection of three threads: confidence-based OSR, DQN-style open-set decision heads, and multi-domain RF feature extraction. The contribution is not a new RF waveform generator, a complete attack-chain planner, or an end-to-end autonomous EW response system. Instead, the paper asks whether a calibrated OSR decision layer can use multi-domain OTA RF preliminary-action evidence to separate known behaviors from novel behaviors under realistic calibration constraints. The resulting known detections can support precursor-level ATT\&CK/EW reasoning, while rejected unknowns can populate a downstream label-making and continual-learning loop.
